\subsection{导子的延拓}

引理 4. 设 V 是 K-模，T 是 V 的张量代数。设 u 是 V 的自同态。存在唯一的延拓 u 的 T 的导子。此导子与 T 上的对称算子交换。

令 \(F = V \times V \times \cdots \times V\)（n 个因子）。映射

\[
\begin{array}{r l} (x _ {1}, \ldots , x _ {n}) & \mapsto u x _ {1} \otimes x _ {2} \otimes \dots \otimes x _ {n} \\ & \qquad + x _ {1} \otimes u x _ {2} \otimes \dots \otimes x _ {n} + \dots + x _ {1} \otimes x _ {2} \otimes \dots \otimes u x _ {n} \end{array}
\]

从 F 到  V 是多线性的。因此存在 \(\bigotimes^{n}\) V 的自同态 \(u_{n}\)\(\bigotimes^{n}\)，使得：

\[
u _ {n} \left(x _ {1} \otimes \dots \otimes x _ {n}\right) = u x _ {1} \otimes \dots \otimes x _ {n} + \dots + x _ {1} \otimes \dots \otimes u x _ {n}
\]

对 V 中所有 \(x_1, \ldots, x_n\) 成立。于是 \(u_1 = u\)。令 v 为 K-模 T 的自同态，它在每个 \(v\) 上与 \(u_n\) 一致，并在 \(\mathrm{T}^n = \bigotimes^n \mathrm{V}\)\(\mathrm{T}^0 = \mathrm{K}. 1\) 上为零。我们证明 v 是 T 的导子。若 \(v\)\(x_1, \ldots, x_n, y_1, \ldots, y_p\) 是 V 的元素，则

\[
\begin{array}{l} v ((x _ {1} \otimes \dots \otimes x _ {n}) \otimes (y _ {1} \otimes \dots \otimes y _ {p})) \\ = \sum_ {i = 1} ^ {n} x _ {1} \otimes \dots \otimes x _ {i - 1} \otimes u x _ {i} \otimes x _ {i + 1} \otimes \dots \otimes x _ {n} \otimes y _ {1} \otimes \dots \otimes y _ {p} \\ \qquad + \sum_ {j = 1} ^ {p} x _ {1} \otimes \dots \otimes x _ {n} \otimes y _ {1} \otimes \dots \otimes y _ {j - 1} \otimes u y _ {j} \otimes y _ {j + 1} \otimes \dots \otimes y _ {p} \\ = v (x _ {1} \otimes \dots \otimes x _ {n}) \otimes (y _ {1} \otimes \dots \otimes y _ {p}) + (x _ {1} \otimes \dots \otimes x _ {n}) \otimes v (y _ {1} \otimes \dots \otimes y _ {p}). \end{array}
\]

\(v\)\(v\)由线性性可知 v 是导子。v 的唯一性是显然的。

最后，显然 \(u_{n}\) 与 \(\otimes V\) 上的所有对称算子交换，因而得到最后一个断言。

命题 7. 设 \(\mathfrak{g}\) 是李代数，U 是其包络代数，\(\sigma\) 是从 \(\mathfrak{g}\) 到 U 的典范映射，D 是 \(\mathfrak{g}\) 的导子。

\begin{enumerate}
\def\labelenumi{(\alph{enumi})}
\item
  存在唯一的 U 的导子 \(\mathrm{D_U}\)，使 \(\sigma \circ \mathrm{D} = \mathrm{D_U} \circ \sigma\)（即当 \(\mathrm{D_U}\)\(\mathfrak{g}\) 在 \(\sigma\) 下可与 U 的一个子模相同时， 延拓 D）。
\item
  \(D_{U}\) 保持 \(U_{n}\) 以及 g 上 n 阶齐次对称张量在 U 中的像的集合 \(U^{n}\) 稳定。
\item
  \(\mathbf{D}_{\mathbf{U}}\) 与 \(\mathbf{U}\) 的主反自同构交换。
\item
  若 \(\mathbf{D}\) 是由 \(\mathfrak{g}\) 的元素 \(x\) 定义的 \(\mathfrak{g}\) 的内导子，则 \(\mathrm{D}_{\mathbf{U}}\) 是由  定义的 \(\mathbf{U}\)\(\sigma(x)\) 的内导子。
\end{enumerate}

令 \(D_T\)\(T\)\(g\)\(D\)\(J\)\(T\) 为延拓 D 的 g 的张量代数 T 的导子（引理 4）。由

\[
x \otimes y - y \otimes x - [ x, y ]
\]

生成的 T 的双边理想 J 在 \((x,y \text{ 属于 } \mathfrak{g})\)\(\mathrm{D}_{\mathrm{T}}\) 下稳定。其中：

\[
\begin{array}{r} \mathrm{D} _ {\mathrm{T}} (x \otimes y - y \otimes x - [ x, y ]) = \mathrm{D} x \otimes y - y \otimes \mathrm{D} x - [ \mathrm{D} x, y ] \\ + x \otimes \mathrm{D} y - \mathrm{D} y \otimes x - [ x, \mathrm{D} y ]. \end{array}
\]

取商后，\(D_{T}\) 定义了 U 的导子 \(D_{U}\)，满足 \(\sigma \circ D = D_{U} \circ \sigma\)。由于 1 和  生成代数 U，\(D_{U}\) 的唯一性是直接的。断言 (b) 是显然的。令 A 为 U 的主反自同构。现在证明 (c)。若 \(\sigma(\mathfrak{g})\)\(x_{1}, \ldots, x_{n}\) 属于 g，则

\[
\begin{array}{r l} \mathrm{D} _ {\mathrm{U}} \mathrm{A} (\sigma (x _ {1}) \dots \sigma (x _ {n})) & = \mathrm{D} _ {\mathrm{U}} ((- 1) ^ {n} \sigma (x _ {n}) \dots \sigma (x _ {1})) \\ & = (- 1) ^ {n} \sum_ {i = 1} ^ {n} \sigma (x _ {n}) \dots \mathrm{D} _ {\mathrm{U}} (\sigma (x _ {i})) \dots \sigma (x _ {1}) \\ & = (- 1) ^ {n} \sum_ {i = 1} ^ {n} \sigma (x _ {n}) \dots \sigma (\mathrm{D} x _ {i}) \dots \sigma (x _ {1}) \\ & = \mathrm{A} \bigg (\sum_ {i = 1} ^ {n} \sigma (x _ {1}) \dots \sigma (\mathrm{D} x _ {i}) \dots \sigma (x _ {n}) \bigg) \\ & = \mathrm{AD} _ {\mathrm{U}} (\sigma (x _ {1}) \dots \sigma (x _ {n})). \end{array}
\]

最后，设 \(x \in \mathfrak{g}\)。令 \(\Delta\) 为 U 的内导子 \(y \mapsto \sigma(x)y - y\sigma(x)\)（Algebra, Chapter IV, § 4, no. 3, Example 2）。则对 \(x' \in \mathfrak{g}\)，

\[
(\Delta \circ \sigma) (x ^ {\prime}) = \sigma (x) \sigma (x ^ {\prime}) - \sigma (x ^ {\prime}) \sigma (x) = \sigma ([ x, x ^ {\prime} ]) = (\sigma \circ \operatorname{ad} x) (x ^ {\prime}),
\]

所以 \(\Delta \circ \sigma = \sigma \circ \operatorname{ad} x\)。证明完毕。

将命题 7 应用于交换李代数的情形可见，K-模的每个自同态 u 都能唯一延拓为该模的对称代数的导子；取商后，此导子由延拓 u 的张量代数导子导出。

再次设 g 是 K 上的李代数，D 是 g 的导子。沿用前面的记号 T、S、U、G。令 \(D_{T}\)、\(D_{S}\) 为分别延拓 D 的 T、S 的导子，令 \(D_{U}\) 为满足 \(\sigma \circ D = D_{U} \circ \sigma\) 的 U 的唯一导子。由于 \(D_{U}\) 保持 \(U_{n}\) 稳定，取商后 \(D_{U}\) 定义了 G 的导子 \(D_{G}\)。由于取商时 \(D_{U}\) 和 \(D_{S}\) 都由 \(D_{T}\) 导出，图 (3) 的交换性表明，\(D_{G}\) 也可由第 6 号定义的同态  从 \(D_{S}\) 导出。若进一步 K 是特征为 0 的域，则图 (4) 的同构将 \(\omega\)\(D_{T}\)、\(D_{S}\)、\(D_{U}\)、\(D_{G}\) 在 \(S^{n}\)、\(S^{n}\)、\(U^{n}\)、\(G^{n}\) 上的限制彼此对应。因此，S 到 U 的典范同构将 \(D_{S}\) 映到 \(D_{U}\)。

